\chapter{Structure du marché européen des actions}\label{ch:m1:european-equity-market-structure}

Une action d'une grande entreprise française se négocie ce matin à la bourse de Paris qui l'a cotée,
sur plusieurs plateformes paneuropéennes, dans les livres propres de quelques
banques, dans des enchères qui durent un dixième
de seconde, et dans des \omterm{def:m1:us-equity-market-structure:dark}{dark pools} qui apparient les ordres à un point médian que personne ne publie
en un seul endroit. En juillet 2026, environ 80~milliards d'euros d'actions européennes ont changé de mains
chaque jour~; à peine la moitié de ce qui s'est négocié sur les plateformes l'a été dans un carnet d'ordres continu,
et visible. Et jusqu'en septembre 2026, le continent n'avait jamais eu ce
que les États-Unis ont depuis longtemps~: un flux public unique indiquant quel est le meilleur
prix. Ce chapitre décrit un marché qui a résolu les mêmes
problèmes que le marché américain, mais par des choix opposés.

\section{Une directive, quatre types de lieux}

La directive sur les marchés d'instruments financiers a mis fin aux
monopoles des bourses nationales en 2007~; sa deuxième version, MiFID~II, accompagnée du règlement
MiFIR, régit les négociations européennes depuis janvier 2018. Elle classe
chaque lieu où une action peut se négocier dans une catégorie juridique.

\begin{definition}[Marché réglementé et système organisé de négociation]\label{def:m1:european-equity-market-structure:rm}
Un \emph{marché réglementé}\index{marché réglementé} est un système multilatéral,
exploité par un opérateur de marché, qui réunit de multiples intérêts acheteurs et vendeurs de tiers
selon des règles non discrétionnaires et admet des instruments à la négociation~: il s'agit de la forme
juridique européenne d'une bourse. Un
\emph{système organisé de négociation}\index{système organisé de négociation} (OTF) est un système multilatéral
dont l'opérateur peut exercer un pouvoir discrétionnaire dans l'appariement des ordres~; il
concerne les obligations et les dérivés et ne peut pas négocier d'actions.
\end{definition}

Entre les deux se situe le \omterm{def:m1:exchanges-brokers-venues:mtf}{système multilatéral de négociation}
(\cref{def:m1:exchanges-brokers-venues:mtf}), qui reprend les règles du \omterm{def:m1:european-equity-market-structure:rm}{marché réglementé} sans sa
fonction de cotation~: les plateformes paneuropéennes sont
des MTF.

\begin{definition}[Internalisateur systématique]\label{def:m1:european-equity-market-structure:si}
Un \emph{internalisateur systématique}\index{internalisateur systématique} (SI) est une
entreprise d'investissement qui, de manière organisée, fréquente et substantielle, négocie
pour son propre compte lorsqu'elle exécute des ordres de clients en dehors d'une plateforme de négociation. Elle
est une contrepartie, et non une plateforme~: chaque transaction est bilatérale, contre le carnet propre
de l'entreprise, et celle-ci doit publier des cotations fermes sur les actions liquides, jusqu'à une
taille standard.
\end{definition}

\begin{omfigure}
\begin{tikzpicture}[font=\small,
  top/.style={draw=omExo,rounded corners=2pt,minimum width=3.4cm,minimum height=0.8cm,align=center,fill=omExo!6},
  box/.style={draw=omDef,rounded corners=2pt,minimum width=2.9cm,minimum height=1.25cm,align=center,fill=omDef!4},
  bil/.style={draw=omThm,rounded corners=2pt,minimum width=2.9cm,minimum height=1.25cm,align=center,fill=omThm!5}]
\node[top] (m) at (4.0,2.2) {Multilatéral\\(plusieurs à plusieurs)};
\node[top] (b) at (12.0,2.2) {Bilatéral\\(entreprise contre client)};
\node[box] (rm) at (0.6,0) {Marché réglementé\\cote et négocie};
\node[box] (mtf) at (4.0,0) {MTF\\négociation seule};
\node[box] (otf) at (7.4,0) {OTF\\hors actions};
\node[bil] (si) at (10.4,0) {Internalisateur\\systématique};
\node[bil] (otc) at (13.6,0) {Autres OTC};
\foreach \x in {rm,mtf,otf}{ \draw[omExo,thick] (m.south) -- (\x.north); }
\foreach \x in {si,otc}{ \draw[omExo,thick] (b.south) -- (\x.north); }
\end{tikzpicture}

\omcaption{Où une action européenne peut légalement se négocier. Une \emph{obligation de négociation des actions} oblige les entreprises
d'investissement à exécuter les actions sur un \omterm{def:m1:european-equity-market-structure:rm}{marché réglementé},
sur un MTF ou auprès d'un \omterm{def:m1:european-equity-market-structure:si}{internalisateur systématique}, ce qui limite le véritable marché de gré à gré
aux cas exceptionnels.}\label{fig:m1:european-equity-market-structure:taxonomy}
\end{omfigure}

\begin{definition}[Meilleure exécution]\label{def:m1:european-equity-market-structure:bestex}
La \emph{meilleure exécution}\index{meilleure exécution} (best execution) est l'obligation
pour une entreprise qui exécute des ordres de clients de prendre toutes les mesures suffisantes pour obtenir le meilleur
résultat possible, en tenant compte du prix, des coûts, de la rapidité, de la probabilité d'exécution
et de règlement-livraison, de la taille et de la nature de l'ordre.
\end{definition}

C'est la réponse européenne à la fragmentation, et elle est à l'opposé de la
réponse américaine. Il n'existe pas de \omterm{def:m1:us-equity-market-structure:opr}{règle de protection des ordres}
(\cref{def:m1:us-equity-market-structure:opr})~: aucune plateforme n'est obligée d'acheminer un ordre
vers un meilleur prix ailleurs, et une transaction à un prix inférieur n'enfreint aucune
règle de marché. Le devoir incombe au \emph{courtier} (broker), sous forme d'une politique qu'il doit rédiger,
appliquer et être en mesure de démontrer. Les plateformes se connectent à qui elles veulent~; les routeurs des
courtiers relient le marché dans son ensemble.

\section{Fragmentation sans flux consolidé}

\begin{definition}[Fournisseur de système consolidé de publication]\label{def:m1:european-equity-market-structure:ctp}
Un \emph{fournisseur de système consolidé de publication}\index{fournisseur de système consolidé de publication} (CTP)
est une entité autorisée à collecter les comptes rendus de transactions et, pour les actions, les meilleurs
cours acheteurs et vendeurs, auprès de toutes les plateformes et de tous les dispositifs de publication, puis
à les diffuser sous la forme d'un flux continu unique.
\end{definition}

\begin{dated}{2026-09}{Un flux enfin disponible}\label{dat:m1:european-equity-market-structure:ctp}
Pendant dix-huit ans après le début de la concurrence, l'Europe n'a eu aucun flux consolidé
pour les actions~: chaque entreprise construisait sa propre vue à partir des flux des plateformes, ou
achetait celle d'un fournisseur. En décembre 2025, l'ESMA a retenu un consortium de bourses, EuroCTP,
comme premier fournisseur pour les actions et les fonds indiciels cotés, puis l'a autorisé
pour cinq ans sous sa supervision directe. Le lancement du flux était fixé
au 14 septembre 2026, avec une période de transition allant jusqu'au 30 septembre
2026. Ses données sont gratuites pour les investisseurs particuliers, les universitaires et les
régulateurs.
\end{dated}

L'absence de meilleur prix officiel a eu des conséquences pratiques qui
perdurent. «~Le meilleur cours acheteur et vendeur européen~» est un chiffre que chaque entreprise calcule
pour elle-même, à partir des plateformes auxquelles elle a choisi de se connecter, sans filtre de \omterm{def:m1:us-equity-market-structure:lots}{lot
standard} et sans statut protégé~; deux courtiers peuvent honnêtement ne pas s'accorder à son sujet.
L'analyse des coûts de transaction doit préciser quelles plateformes son indice de référence incluait.
Et le prix de la bourse principale conserve, par habitude, un rôle privilégié~: lorsque la bourse
nationale suspend ses cotations, les volumes sur les plateformes qui recopient son prix se tarissent
souvent aussi.

\begin{definition}[Nombre effectif de plateformes]\label{def:m1:european-equity-market-structure:hhi}
Si les plateformes ont des parts de marché $s_1,\dots,s_n$ dont la somme vaut un,
l'\emph{indice de Herfindahl} vaut $H = \sum_i s_i^2$ et le \emph{nombre effectif de plateformes}
est $1/H$~: c'est le nombre de plateformes de même taille qui produirait la
même concentration.
\end{definition}

\begin{example}[Fragmenté, mais dans quelle mesure~?]\label{ex:m1:european-equity-market-structure:hhi}
Une action se négocie à 60~\% sur sa bourse principale, à 20~\% sur une plateforme et à 10~\% sur
chacune de deux autres. $H = 0.36 + 0.04 + 0.01 + 0.01 = 0.42$~: soit un nombre effectif de 2.4
plateformes, alors qu'il en existe quatre. Si la part de la bourse principale tombe à 40~\%
et que les autres en prennent 20~\% chacune, $H = 0.28$ et le nombre effectif vaut 3.6.
La mesure dépend surtout de la part du leader
(\cref{fig:m1:european-equity-market-structure:effective}).
\end{example}

\begin{omfigure}
\begin{tikzpicture}
\begin{axis}[width=10.5cm,height=4.6cm,
  xlabel={part de marché de la plus grande plateforme (\%)},ylabel={nombre effectif de plateformes},
  tick label style={font=\small},label style={font=\small},
  xmin=20,xmax=100,ymin=1,ymax=6,grid=major,grid style={gray!25},table/col sep=comma]
\addplot[omDef,very thick] table[x=leader_pct,y=effective]{figdata/markets-1/11-european-equity-market-structure/effective.csv};
\end{axis}
\end{tikzpicture}

\omcaption{Nombre effectif de plateformes lorsqu'une plateforme détient la part figurant en
abscisse et que cinq autres se partagent le reste à parts égales. Six plateformes ne valent
«~six~» que si le leader ne représente plus qu'un sixième du marché. Données~: calcul
effectué par le script du chapitre.}\label{fig:m1:european-equity-market-structure:effective}
\end{omfigure}

\section{Les mécanismes~: cotés, sombres, périodiques, clôture}

\begin{dated}{2026-09}{Comment les actions européennes se sont négociées en juillet 2026}\label{dat:m1:european-equity-market-structure:july}
Selon les statistiques mensuelles d'un opérateur paneuropéen, le volume échangé
en bourse sur les actions européennes s'est élevé en moyenne à 57.5~milliards d'euros
par jour en juillet 2026, sur 80~milliards d'euros d'activité accessible, hors bourse
compris. Dans la part réalisée en bourse, les carnets d'ordres continus à cours limité
ont représenté 52.8~\%, les enchères de clôture 24.3~\%, les \omterm{def:m1:european-equity-market-structure:periodic}{enchères périodiques}
10.4~\% et les carnets d'ordres non affichés 10.9~\%. Les \omterm{def:m1:european-equity-market-structure:si}{internalisateurs systématiques}
ont traité 14.3~milliards d'euros par jour, soit 17.8~\% du volume accessible.
\end{dated}

\begin{omfigure}
\begin{tikzpicture}
\begin{axis}[width=9cm,height=4.2cm,xbar,bar width=8pt,
  xlabel={part du volume échangé en bourse, juillet 2026 (\%)},
  ytick=data,yticklabels={{Carnets d'ordres centraux},{Enchères de clôture},{Enchères périodiques},{Carnets d'ordres non affichés}},
  yticklabel style={font=\small},tick label style={font=\small},label style={font=\small},
  y dir=reverse,xmin=0,xmax=65,enlarge y limits=0.2,grid=major,grid style={gray!25},
  nodes near coords,nodes near coords style={font=\footnotesize},
  table/col sep=comma]
\addplot[fill=omDef!60,draw=omDef] table[x=pct,y=k]{figdata/markets-1/11-european-equity-market-structure/mechanisms.csv};
\end{axis}
\end{tikzpicture}

\omcaption{Quatre modes de confrontation des ordres sur les plateformes européennes. Un quart
du volume échangé en bourse se réalise pendant les quelques minutes des enchères de clôture.
Données~: comme dans l'\cref{dat:m1:european-equity-market-structure:july}.}
\end{omfigure}

MiFID II impose la transparence avant négociation~: une plateforme doit publier ses
cours acheteur et vendeur. La négociation sombre existe par \emph{dérogation} à cette règle.

\begin{definition}[Dérogation au prix de référence et dérogation pour taille élevée]\label{def:m1:european-equity-market-structure:waivers}
Au titre de la \emph{dérogation au prix de référence}\index{dérogation au prix de référence}, une plateforme
n'est pas tenue d'afficher les ordres exécutés à un prix repris d'un autre marché, en pratique le point médian
de la cotation de la bourse principale. Au titre de la
\emph{dérogation pour taille élevée}\index{dérogation pour taille élevée}, elle n'est pas tenue d'afficher un ordre
qui est grand au regard de la taille normale de marché pour cette action.
\end{definition}

\begin{definition}[Plafond de volume]\label{def:m1:european-equity-market-structure:cap}
Le \emph{double plafond de volume}\index{double plafond de volume} initial suspendait la négociation
d'une action au titre de la \omterm{def:m1:european-equity-market-structure:waivers}{dérogation au prix de référence} lorsque, sur douze mois, ce type de négociation
dépassait 4~\% du volume total de l'action sur une plateforme quelconque ou 8~\% à l'échelle de l'Union.
Depuis octobre 2025, il est remplacé par un \emph{plafond de volume
unique}~: 7~\% à l'échelle de l'Union, avec une suspension de trois mois de la dérogation pour
cette action lorsque ce seuil est dépassé.
\end{definition}

\begin{dated}{2026-09}{Le plafond en pratique}\label{dat:m1:european-equity-market-structure:svc}
L'ESMA a publié les premiers résultats du plafond de volume unique le 9 octobre 2025 et publie
les calculs et les suspensions chaque trimestre~; le système du \omterm{def:m1:european-equity-market-structure:cap}{double plafond de volume} devait
être démantelé en janvier 2026. La négociation en taille élevée
échappe au plafond.
\end{dated}

\begin{definition}[Enchère périodique]\label{def:m1:european-equity-market-structure:periodic}
Une \emph{enchère périodique}\index{enchère périodique} est une courte enchère, qui dure typiquement
une fraction de seconde, et qui débute lorsque deux ordres du carnet pourraient se croiser,
publie un cours indicatif et une taille indicative pendant sa brève existence, puis se conclut
à un prix unique, souvent contraint de se situer entre
le meilleur cours acheteur et le meilleur cours vendeur du marché de référence.
\end{definition}

Les \omterm{def:m1:european-equity-market-structure:periodic}{enchères périodiques} se sont développées lorsque les premiers plafonds ont suspendu la négociation sombre pour beaucoup
d'actions~: une action suspendue de la négociation sombre au point médian pouvait encore se traiter au point médian, ou
près de lui, dans une enchère «~cotée~» à la lettre de la loi, qui
livre presque aussi peu d'information qu'un \omterm{def:m1:us-equity-market-structure:dark}{dark pool}. Chaque règle de ce marché
a produit le mécanisme qui la contourne~; le lecteur qui veut prévoir la
suivante doit se demander ce que le nouveau plafond unique rend coûteux.

\begin{omfigure}
\begin{tikzpicture}
\begin{axis}[width=11cm,height=5cm,
  xlabel={mois},ylabel={part du volume de l'action (\%)},
  tick label style={font=\small},label style={font=\small},
  xmin=12,xmax=36,ymin=0,ymax=14,grid=major,grid style={gray!25},unbounded coords=jump,
  legend columns=1,legend cell align=left,legend style={font=\footnotesize,at={(0.5,-0.32)},anchor=north,draw=none,fill=none,/tikz/every even column/.append style={column sep=6pt}},
  table/col sep=comma]
\addplot[omThm,very thick] table[x=month,y=usage_pct]{figdata/markets-1/11-european-equity-market-structure/cap.csv};
\addplot[omExo,thick,const plot mark mid] table[x=month,y=dark_pct]{figdata/markets-1/11-european-equity-market-structure/cap.csv};
\addplot[omDef,thick,const plot mark mid] table[x=month,y=periodic_pct]{figdata/markets-1/11-european-equity-market-structure/cap.csv};
\addplot[black,dashed] coordinates {(12,7) (36,7)};
\legend{utilisation de la dérogation sur douze mois,part mensuelle de la négociation sombre,part mensuelle des enchères périodiques}
\end{axis}
\end{tikzpicture}

\omcaption{Une action simulée dont la part de négociation sombre au point médian dérive à la hausse. Lorsque
l'utilisation sur douze mois franchit 7~\% (tirets), la dérogation est suspendue pendant
trois mois~: le volume sombre tombe à zéro et la moitié de celui-ci réapparaît dans les
\omterm{def:m1:european-equity-market-structure:periodic}{enchères périodiques}. Données~: simulation du chapitre~; la répartition des migrations
est illustrative.}\label{fig:m1:european-equity-market-structure:cap}
\end{omfigure}

\section{La clôture et le tick}

Un quart du chiffre d'affaires sur bourse dans l'enchère de clôture n'est pas un hasard
de goût. Les fonds indiciels doivent négocier au cours de clôture auquel ils sont
comparés (\cref{ch:m1:the-buy-side})~; la bourse de rattachement détient ce prix~;
et plus le volume s'y concentre, plus le carnet d'ordres intrajournalier s'amincit,
ce qui attire encore davantage de volume. Le mécanisme d'enchère lui-même fait l'objet du
\cref{ch:m1:auctions}.

Contrairement à l'incrément américain d'un cent, les pas de cotation européens sont fixés par
un tableau. Dans le régime des pas de cotation (RTS~11), l'incrément minimal d'une action
dépend de deux paramètres~: son cours, et une \emph{bande de liquidité} donnée par son
nombre moyen quotidien de transactions sur son marché le plus pertinent, recalculé une fois
par an. Les actions plus liquides ont des ticks plus petits à cours égal. La conception
vise une fourchette de quelques ticks partout~; elle implique aussi que le tick d'une action peut
changer le premier lundi d'avril sans aucune modification de son cours, ce qui modifie
la longueur des files d'attente et l'économie des \omterm{def:m1:what-a-trading-firm-does:market-maker}{teneurs de marché} du jour au lendemain.

\section{Tutoriel~: mesurer la fragmentation}\label{tut:m1:european-equity-market-structure:frag}

\begin{tutorial}
\textbf{Objectif.} Construire le meilleur cours acheteur et vendeur européen, mesurer la
concentration et surveiller un plafond de volume. \textbf{État final~:}
\cref{fig:m1:european-equity-market-structure:effective,fig:m1:european-equity-market-structure:cap}.
\begin{tutsteps}
\item \textbf{Concentration et meilleure cotation.} Il n'y a ni filtre de \omterm{def:m1:us-equity-market-structure:lots}{lot standard}
ni statut protégé~: le meilleur prix est simplement le meilleur prix parmi les
plateformes retenues pour l'examen.
\omcode{code/markets-1/11-european-equity-market-structure/python/euro_frag.py}{11}{27}{Indice de Herfindahl, nombre effectif de plateformes et meilleur cours acheteur et vendeur entre plateformes.}\label{lst:m1:european-equity-market-structure:hhi}
\item \textbf{Essayer.} Avec Paris à $45.20 \times 45.24$ et deux plateformes
à $45.21 \times 45.23$ et $45.21 \times 45.24$, la meilleure cotation est $45.21
\times 45.23$~: plus serrée que celle de la bourse de rattachement, et invisible pour
quiconque ne regarde que la bourse de rattachement.
\item \textbf{Le plafond est un ratio glissant.}
\omcode{code/markets-1/11-european-equity-market-structure/python/euro_frag.py}{30}{34}{Utilisation sur douze mois de la dérogation au prix de référence.}\label{lst:m1:european-equity-market-structure:cap}
\item \textbf{Simuler un dépassement.} Dans le scénario du chapitre, l'utilisation franchit
7~\% au mois~28~; le carnet sombre est vide aux mois~29 à 31 et la part des \omterm{def:m1:european-equity-market-structure:periodic}{enchères
périodiques} passe de 8~\% à environ 12~\%.
\end{tutsteps}
\textbf{Pour aller plus loin.} Faire de l'utilisation une prévision~: étant donné les onze
derniers mois, quelle part sombre le mois suivant déclenche une suspension~? Le routeur d'un
courtier a besoin exactement de ce nombre. Recalculer ensuite
l'\cref{ex:m1:european-equity-market-structure:hhi} en comptant les \omterm{def:m1:european-equity-market-structure:si}{internalisateurs
systématiques} comme une «~plateforme~» de plus.
\end{tutorial}

\section{Construction~: le moniteur de fragmentation}\label{bld:m1:european-equity-market-structure:monitor}

\begin{build}
\bfield{Rôle} Le routeur de l'entreprise fictive et son analyse de coûts ont tous deux besoin
de savoir, par instrument, où le volume se négocie et selon quel mécanisme, et à quel point
chaque action est proche de perdre sa liquidité sombre.
\bfield{Interface} \texttt{Monitor.add(trade)} avec une transaction portant
l'instrument, le code de plateforme (\cref{bld:m1:exchanges-brokers-venues:registry}),
le mécanisme (\texttt{lit}, \texttt{dark\_rpw}, \texttt{dark\_lis},
\texttt{periodic}, \texttt{auction}, \texttt{si}), la valeur et le mois~;
\texttt{shares(instrument, by)} avec \texttt{by} égal à \texttt{venue} ou
\texttt{mechanism}; \texttt{effective\_venues(instrument)};
\texttt{cap\_usage(instrument, month)}; \texttt{headroom(instrument, month)}
qui renvoie le volume de dérogation qui ramènerait l'utilisation du mois suivant au plafond.
\bfield{Règles} Les valeurs sont des entiers en unités du registre. Seul \texttt{dark\_rpw}
compte pour le plafond~; la \omterm{def:m1:european-equity-market-structure:waivers}{dérogation pour taille élevée} n'y compte pas. L'utilisation exige douze mois
complets~: avant cela, le moniteur ne renvoie aucun chiffre, et non un chiffre partiel.
\bfield{Tests d'acceptation} \texttt{code/firm/fragmentation/tests/}~: les parts somment à
un~; l'indice de Herfindahl de l'exemple~; l'utilisation du plafond sur une série constante~;
l'exclusion de la \omterm{def:m1:european-equity-market-structure:waivers}{dérogation pour taille élevée}~; le headroom ramène exactement l'utilisation à 7~\%.
\bfield{Pour aller plus loin} Un calendrier de suspensions~: étant donné les dates de publication
trimestrielles, les mois pendant lesquels la dérogation de chaque instrument est indisponible.
\end{build}

\begin{omsources}
\item Directive 2014/65/UE (MiFID II), article 4(1)(20)--(24) et article 27~;
règlement (UE) n°~600/2014 (MiFIR) tel que modifié en 2024.
\item ESMA, \emph{ESMA selects EuroCTP to become the first Consolidated Tape
Provider for shares and ETFs}, 19 décembre 2025~; \emph{ESMA authorises
EuroCTP as the Consolidated Tape Provider for shares and exchange-traded
funds}, 2026.
\item ESMA, \emph{ESMA prepares for switch toward single volume cap in October
2025}~; \emph{Double Volume Cap Mechanism} (abandonné).
\item Règlement délégué (UE) 2017/588 de la Commission (RTS 11), régime des pas de cotation
pour les actions, les certificats de dépôt et les fonds indiciels cotés.
\item Cboe Global Markets, \emph{Market Metrics That Matter: European Equities
July Volume Briefing}, 17 août 2026.
\item C.-A.~Lehalle et S.~Laruelle (éd.), \emph{Market Microstructure in
Practice}, 2e éd., World Scientific, 2018, chapitre~1.
\end{omsources}

\section{Exercices}

\begin{exercise}[$\star$]\label{exo:m1:european-equity-market-structure:1}
Classer chacun comme \omterm{def:m1:european-equity-market-structure:rm}{marché réglementé}, MTF, OTF, \omterm{def:m1:european-equity-market-structure:si}{internalisateur systématique} ou
autre~: (a) la bourse sur laquelle une entreprise allemande est cotée~; (b) une
plateforme paneuropéenne négociant 3\,000 actions qu'elle n'a pas cotées~; (c) une banque
exécutant toute la journée les ordres sur actions de ses clients depuis son propre livre~;
(d) le système d'un courtier appariant des ordres d'obligations à discrétion.
\end{exercise}

\begin{exercise}[$\star$]\label{exo:m1:european-equity-market-structure:2}
Le volume d'une action se répartit à 45~\%, 25~\%, 15~\%, 10~\% et 5~\% entre cinq plateformes.
Calculer l'indice de Herfindahl et le nombre effectif de plateformes.
\end{exercise}

\begin{exercise}[$\star$]\label{exo:m1:european-equity-market-structure:3}
À partir de l'\cref{dat:m1:european-equity-market-structure:july}, calculer en euros par
jour le chiffre d'affaires des enchères de clôture, des \omterm{def:m1:european-equity-market-structure:periodic}{enchères périodiques} et des
carnets non affichés, ainsi que l'activité hors bourse qui n'est pas celle des
\omterm{def:m1:european-equity-market-structure:si}{internalisateurs systématiques}.
\end{exercise}

\begin{exercise}[$\star\star$]\label{exo:m1:european-equity-market-structure:4}
Trois plateformes cotent une action~: X $18.340 \times 18.360$, Y $18.345 \times
18.355$, Z $18.350 \times 18.365$. Donner le meilleur cours acheteur et le meilleur cours vendeur, ainsi que leurs plateformes.
Un courtier connecté uniquement à X et Y achète à 18.355. A-t-il enfreint une
règle de marché~? A-t-il failli envers son client~?
\end{exercise}

\begin{exercise}[$\star\star$]\label{exo:m1:european-equity-market-structure:5}
Au cours des onze derniers mois, une action s'est échangée pour un total de 4\,400 millions d'euros, dont
300 millions d'euros sous la \omterm{def:m1:european-equity-market-structure:waivers}{dérogation au prix de référence}. Le total du mois prochain
devrait atteindre 400 millions d'euros. Quel volume sous dérogation le mois prochain porte exactement l'utilisation sur
douze mois à 7~\%~? L'exprimer en part du volume de ce
mois.
\end{exercise}

\begin{exercise}[$\star\star$]\label{exo:m1:european-equity-market-structure:6}
Une action cotée 12.40~\euro a un tick de 0.005~\euro et un spread typique de deux
ticks. Après le recalcul annuel, sa bande de liquidité change et le tick
passe à 0.01~\euro. Donner le spread en points de base avant, et le spread minimal
possible après. Que se passe-t-il pour la file d'attente au meilleur prix~?
\end{exercise}

\begin{exercise}[$\star\star\star$]\label{exo:m1:european-equity-market-structure:7}
\emph{Codage.} Avec \texttt{simulate\_cap}, trouver le premier mois de dépassement pour
les graines 0 à 49 et en donner la médiane. Mettre ensuite à zéro la dérive haussière de la
part sombre et indiquer combien des cinquante trajectoires dépassent encore le plafond.
\end{exercise}

\begin{exercise}[$\star\star\star$]\label{exo:m1:european-equity-market-structure:8}
\emph{Trouver la faille.} Un rapport de coûts de transaction affirme~: «~Notre algorithme
a obtenu une exécution moyenne de 0.4 point de base meilleure que le meilleur cours
acheteur et vendeur européen.~» L'indice de référence est construit à partir des trois plateformes sur lesquelles
l'algorithme opère~; l'action se négocie aussi sur quatre autres plateformes et avec sept
\omterm{def:m1:european-equity-market-structure:si}{internalisateurs systématiques}. Expliquer ce que le chiffre montre et ne montre pas, et
comment le corriger.
\end{exercise}

\section{Problème~: où est passé le volume~?}

\begin{problem}[{Problème du week-end --- une action perd sa dark pool}]\label{pb:m1:european-equity-market-structure:1}
Une action à forte capitalisation se négocie en moyenne pour 500 millions d'euros par jour, 21
jours par mois. Son chiffre d'affaires se répartit comme suit~: le carnet continu de la
bourse principale 38~\%, l'enchère de clôture de la principale 20~\%, les carnets affichés de trois MTF à 9~\% chacun (27~\%),
carnets noirs à cours médian sous la \omterm{def:m1:european-equity-market-structure:waivers}{dérogation au prix de référence} 7.5~\%, \omterm{def:m1:european-equity-market-structure:periodic}{enchères périodiques}
4~\%, blocs pour taille élevée 3.5~\%. (Les \omterm{def:m1:european-equity-market-structure:si}{internalisateurs systématiques} sont ignorés
ici.)

\medskip\noindent\textbf{Partie I --- Avant.}
\begin{enumerate}
\item Donner le chiffre d'affaires mensuel et la valeur mensuelle échangée dans les
carnets noirs à cours médian.
\item Par \emph{opérateur de plateforme}, les parts sont les suivantes~: bourse principale 60~\% (son carnet, son
enchère et 2~\% du volume sombre et périodique), le reste se répartit à parts égales
entre quatre opérateurs de plateforme. Calculer le nombre effectif d'opérateurs.
\item Par \emph{mécanisme}, calculer l'indice de Herfindahl sur les six
catégories ci-dessus, puis le nombre effectif de mécanismes.
\item L'utilisation de la dérogation sur douze mois s'établit à 7.5~\%. Que se passe-t-il, et
pendant combien de temps~?
\end{enumerate}

\medskip\noindent\textbf{Partie II --- La suspension.}
Pendant la suspension, les 7.5~\% qui se négociaient dans les carnets noirs à cours médian
se répartissent ainsi~: 45~\% vont aux \omterm{def:m1:european-equity-market-structure:periodic}{enchères périodiques}, 25~\% aux blocs pour taille élevée, 20~\%
aux carnets continus affichés (répartis entre la bourse principale et les MTF au prorata de leurs parts affichées),
et 10~\% ne sont pas négociés du tout.
\begin{enumerate}[resume]
\item Donner la nouvelle part de chaque mécanisme, en pourcentage du
chiffre d'affaires \emph{initial}.
\item Renormaliser par rapport au nouveau chiffre d'affaires, plus faible, et donner les nouveaux pourcentages.
\item De combien de points de pourcentage les \omterm{def:m1:european-equity-market-structure:periodic}{enchères périodiques} ont-elles gagné~?
\item Recalculer le nombre effectif de mécanismes.
\item Les exécutions à cours médian d'un courtier institutionnel lui ont permis d'économiser la moitié du spread, soit 2
points de base, sur le volume qu'il a négocié dans les carnets noirs, lequel représentait 12~\% de ses 40 millions d'euros
par jour sur cette action. Si seules les fractions de ce flux correspondant aux \omterm{def:m1:european-equity-market-structure:periodic}{enchères périodiques} et aux blocs
continuent de s'exécuter à cours médian, combien la suspension lui coûte-t-elle
par mois~?
\end{enumerate}

\medskip\noindent\textbf{Partie III --- Le retour.}
\begin{enumerate}[resume]
\item Après trois mois, la dérogation est rétablie. Avec un volume sombre nul pendant
trois mois et 7.5~\% par mois auparavant, quelle est l'utilisation sur douze mois
le jour du rétablissement~?
\item Si la négociation sombre reprend à 7.5~\% par mois, combien de mois faut-il
pour que l'utilisation dépasse à nouveau 7~\%~?
\item Quelle part sombre mensuelle maintient l'utilisation à exactement 7~\% indéfiniment~?
\item Pourquoi les opérateurs de carnets noirs pourraient-ils collectivement préférer rester en deçà
de ce niveau, et pourquoi ne peuvent-ils pas facilement se coordonner pour le faire~?
\end{enumerate}

\medskip\noindent\textbf{Partie IV --- Jugement.}
\begin{enumerate}[resume]
\item Qui est protégé par le plafond, selon le raisonnement du législateur~?
\item Qui en supporte le coût, selon la question 9~?
\item Les \omterm{def:m1:european-equity-market-structure:periodic}{enchères périodiques} affichent un cours indicatif pendant une fraction de
seconde. En quel sens sont-elles cotées, et en quel sens sont-elles sombres~?
\item Pourquoi la \omterm{def:m1:european-equity-market-structure:waivers}{dérogation pour taille élevée} échappe-t-elle au plafond~?
\item Un système consolidé de publication diffusant désormais le meilleur cours acheteur et vendeur, quelles
statistiques de ce problème deviennent plus faciles à calculer, et lesquelles
ne changent pas~?
\item Énoncer le \emph{résultat nommé}~: la part de marché des \omterm{def:m1:european-equity-market-structure:periodic}{enchères
périodiques} pendant la suspension, en pourcentage du chiffre d'affaires renormalisé.
\item En une phrase~: quelle est la loi générale de la structure de marché
qu'illustre ce problème~?
\end{enumerate}
\end{problem}

\section{Questions d'entretien}

\begin{interviewq}[$\star$ \iqroles{trader, chercheur, développeur}]\label{iq:m1:european-equity-market-structure:1}
Citez les principales différences entre la structure des marchés d'actions américains et européens.
\end{interviewq}

\begin{interviewq}[$\star$ \iqroles{trader, banque}]\label{iq:m1:european-equity-market-structure:2}
Qu'est-ce qu'un \omterm{def:m1:european-equity-market-structure:si}{internalisateur systématique}, et en quoi diffère-t-il d'un
\omterm{def:m1:exchanges-brokers-venues:mtf}{système multilatéral de négociation}~?
\end{interviewq}

\begin{interviewq}[$\star\star$ \iqroles{chercheur, trader}]\label{iq:m1:european-equity-market-structure:3}
Pourquoi un quart du volume européen négocié en bourse est-il traité dans
l'enchère de clôture, et qu'est-ce que cela fait à la liquidité intrajournalière~?
\end{interviewq}

\begin{interviewq}[$\star\star$ \iqroles{développeur, chercheur}]\label{iq:m1:european-equity-market-structure:4}
Vous devez calculer un meilleur cours acheteur et vendeur européen pour l'analyse des coûts de transaction.
Quelles décisions devez-vous prendre que n'aurait pas à prendre un analyste américain~?
\end{interviewq}

\begin{interviewq}[$\star\star$ \iqroles{trader, chercheur}]\label{iq:m1:european-equity-market-structure:5}
Une action est sur le point d'être suspendue de la négociation sombre à cours médian en vertu du plafond de
volume. Comment cela modifie-t-il votre façon d'exécuter un ordre important sur celle-ci~?
\end{interviewq}

\begin{interviewq}[$\star\star\star$ \iqroles{chercheur, trader}]\label{iq:m1:european-equity-market-structure:6}
Le pas de cotation d'une action double après le recalcul annuel. Prévoyez
l'effet sur le spread, la profondeur, les délais en file d'attente et la rentabilité d'un \omterm{def:m1:what-a-trading-firm-does:market-maker}{teneur de marché}, puis
indiquez comment vous testeriez vos prévisions.
\end{interviewq}
